\scanpage{121}
$\mathcal X$ sous $\unsure{\gamma|\Gamma}$, et de la transitivité de $\unsure{\gamma|\Gamma}$ sur $J$, on voit que cet idéal $\unsure{\mathfrak a}_j$ ne dépend pas de $j$. D'autre part, l'intersection des $\unsure{\mathfrak a}_j$ est évidemment le noyau de
\[
 \deleted{\unsure{k}}k\longrightarrow\Gamma(\unsure{X_J},\mathcal O_{\unsure{X_J}})\simeq k^J.
\]
Comme celui-ci est nul, on a $\unsure{\mathfrak a}_j=0$ \unsure{pour|puis} $\mathcal X_j=X_j$ donc $\mathcal X=X_J$ O.K.
\figure{\begin{tikzpicture}\draw (0,.15) .. controls (.05,.03) and (.25,0) .. (.3,-.1);\end{tikzpicture}}

Ce résultat vaut sans restriction de caractéristique, mais il concerne la situation doublement étale, \uline{trop} restrictive. J'ai envie de vérifier que dès qu'on a trois rev. quadratiques $X_1,X_2,X_3$ de $X$, et une section $\varepsilon$ de $X_1\wedge X_2\wedge X_3$, sans condition étale, on en déduit sur
\[
 \unsure{X|\mathcal X}=\pi^{-1}(\varepsilon)
\]
\margin{\diagram{\begin{tikzcd}[column sep=large]
\unsure{\subset}X_1\times_X X_2\times_X X_3\arrow[r,"\pi"]&X_1\wedge X_2\wedge X_3
\end{tikzcd}}}
une structure biquadratique \unsure{p.r.} aux $\unsure{X_i}$. Ceci doit généraliser la structure biquadratique can. sur un $X_1\times X_2$, par les applications
\diagram{\begin{tikzcd}[column sep=large]
&X_1\\
X_1\times X_2\arrow[ur]\arrow[r,no head]\arrow[dr]&X_2\\
&X_3=X_1\wedge X_2
\end{tikzcd}},
comme on va le voir \unsure{de fait|de suite}.
\mspage{120}
